\section{Bölüm sonu çözümlü problemler}

\begin{enumerate}
  \item \textbf{Araştırma sorusunu yapılandırma.} Bir danışma merkezi, ``akran desteğinin öğrencilerin okula uyumunu artırıp artırmadığını'' incelemek istiyor. Hedef evreni, yanıtı, açıklayıcı değişkeni ve uygun temel tasarımı belirtin.

  \textbf{Çözüm:} Hedef evren ilgili kurumun programa uygun öğrencileridir. Yanıt, geçerliği desteklenmiş bir uyum ölçeğinin son-test puanı; açıklayıcı değişken program grubudur. En güçlü temel tasarım, başlangıç ölçümü bulunan ve öğrencilerin program ile karşılaştırma koşullarına rastgele atandığı ön-test--son-test tasarımıdır. Rastgele atama yapılamıyorsa sonuç ilişki diliyle sunulmalı ve başlangıç farklılıkları incelenmelidir.

  \item \textbf{Parametre ve istatistik.} Üniversitedeki 4.000 birinci sınıf öğrencisinin yalnızlık ortalaması araştırılmaktadır. Rastgele seçilen 160 öğrencinin ortalaması 42,8 bulunmuştur. Parametre ve istatistiği ayırın.

  \textbf{Çözüm:} Parametre, 4.000 öğrencinin bilinmeyen evren ortalaması $\mu$'dür. Örneklemden hesaplanan $\bar{x}=42{,}8$ ise bu parametrenin istatistiksel tahminidir. 160 sayısı örneklem hacmi, 4.000 sayısı evren büyüklüğüdür.

  \item \textbf{Nedensel yorum sınırı.} Programa gönüllü katılan öğrencilerin kaygı ortalaması katılmayanlardan düşük bulunmuştur. ``Program kaygıyı azalttı'' sonucu neden doğrudan kurulamaz?

  \textbf{Çözüm:} Gönüllü katılım seçim yanlılığı yaratabilir. Motivasyon, başlangıç kaygısı, zaman uygunluğu veya yardım arama eğilimi hem katılımla hem sonuçla ilişkili olabilir. Zaman sırası, uygun karşılaştırma, atama süreci ve kayıplar değerlendirilmeden gözlenen fark nedensel etki olarak yorumlanamaz.
\end{enumerate}